\ifx\danslelivre\undefined
\documentclass[10pt]{../fichiers-configuration-latex/smfart}
\usepackage{ifthen}
\newboolean{english}
\newboolean{francais}
\setboolean{english}{false}
\setboolean{francais}{true}
\usepackage{palatino,euler}
\input{../fichiers-configuration-latex/paquets}
\selectlanguage{francais}
\input{../fichiers-configuration-latex/lemme-proposition-theoreme}
\input{../fichiers-configuration-latex/numerotation}
\input{../fichiers-configuration-latex/macros}
\input{../fichiers-configuration-latex/ucs_manquants}
%\input{../.cv}

\usepackage{tikz}
\usetikzlibrary{matrix}
\usetikzlibrary{calc}
\usetikzlibrary{decorations.markings}

\externaldocument{../00-introduction/introduction}
\externaldocument{../02-topologies/topologies}
\externaldocument{../03-approximation/approximation}
\externaldocument{../04-Cohen-Gabber/Cohen-Gabber}
\externaldocument{../07-uniformisation-locale/uniformisation-faible}
\externaldocument{../12-descente-cohomologique-orientee/descente-orientee}
\externaldocument{../13-theoreme-de-finitude/finitude}
\externaldocument{../14-fonctions-dimension/fonctions-dimension}
\externaldocument{../16-classes-de-Chern/gysin}

\begin{document}

\begin{center}
XV. Théorème de Lefschetz affine\\
Vincent Pilloni et Benoît Stroh
\end{center}
%\version

\else
\setboolean{english}{false}
\setboolean{francais}{true}
\selectlanguage{francais}
\chapter{Théorème de Lefschetz affine}
\addtocontents{toc}{par Vincent Pilloni et Benoît Stroh}
\begin{center}
Vincent Pilloni et Benoît Stroh
\end{center}
\fi

\section{Énoncé du théorème et premières réductions}

\subsection{Énoncé}
\label{15-énoncés}

\subsubsection{}
Soient $X$ un schéma muni d'une fonction de dimension $\delta_X$
(\refext{XIV}{ddefi4}) et $n$ un entier inversible sur $X$.
Pour tout faisceau étale $ℱ$ de $𝐙/n$-modules sur $X$,
$$\delta_X(\mc F) = \mathrm{sup}\: \big\{ \delta_X(x) \ , \ x\in X \ | \
{\mc F}_{\xb} \neq 0 \big\}.$$
Rappelons (\refext{XIV}{14-coro-fonction-dim-induite-tf})
qu'un morphisme de type fini $f: Y → X$ induit une
fonction de dimension sur $Y$ ; nous la noterons ici $f^{⋆}δ_X$.
Le théorème principal de cet exposé est le suivant (voir aussi
\refext{0}{0-Lefschetz affine}).

\begin{theoreme2}
\label{15-theoreme-principal}
Supposons le schéma $X$ quasi-excellent et le morphisme $f:Y \ra X$ 
\emph{affine} de type fini. Alors, pour tout faisceau constructible $ℱ$
de $𝐙/n$-modules sur $Y$, on a :
\[
\delta_X(\R^{q}f_⋆(\mc{F})) ≤  f^⋆\delta_X(\mc{F}) - q .
\]
\end{theoreme2}

\begin{remarque2}
\label{15-LA sur trait}
Ce théorème a été déjà démontré en 1994 par O. Gabber
lorsque $X$ est de type fini sur un trait,
cf. \cite{PV@Illusie}. La démonstration du théorème
précédent procède notamment par réduction à ce cas.
\end{remarque2}

\subsection{Reformulation et réductions}

\subsubsection{}
Soient $f$ et $ℱ$ comme ci-dessus.
La conclusion du théorème signifie que pour
tout point géométrique $\xb$ de $X$ localisé en un
point $x$ et tout entier $q>f^⋆δ_X(ℱ)-δ_X(x)$, on a
\[
(\R f_⋆ ℱ)_{\xb}=H^q(Y_{(\xb)},ℱ)=0,
\]
où l'on note $Y_{(\xb)}$ le produit fibré $Y×_X X_{(\xb)}$.
Rappelons (\refext{XIV}{14-coro-schema-tf-sur-corps-fonc-dim})
que le schéma strictement local $X_{(\xb)}$
peut être muni la fonction de dimension $δ_{X_{(\xb)}}:t↦ \dim(\sur{\{t\}})$
(\refext{XIV}{14-coro-schema-tf-sur-corps-fonc-dim}) ;
c'est l'unique fonction de dimension nulle en $x$.
Notons l'inégalité $f^⋆δ_X(ℱ)-δ_X(x) ≥ f_{(\xb)}^⋆δ_{X_{(\xb)}}(ℱ)$,
triviale dans le cas où $δ_X(x)=0$, auquel on peut se ramener.
Ainsi, le théorème \ref{15-theoreme-principal} est équivalent
à l'énoncé suivant.

\begin{corollaire2}
\label{15-prop-principale-locale}
Supposons de plus le schéma $X$ quasi-excellent, \emph{strictement local},
muni de la fonction de dimension $δ_X:t↦ \dim(\sur{\{t\}})$.
Alors, pour tout faisceau constructible $ℱ$ de $𝐙/n$-modules
sur $Y$, on a :
\[
\HH^q(Y,\mc{F})=0 \text{ si } q > f^⋆\delta_X(\mc{F}).
\]
\end{corollaire2}

\subsubsection{}Procédant comme en \sga{4}{xiv}{4.4} pour se ramener au cas d'une
immersion ouverte affine puis utilisant la méthode de la trace
(\sga{4}{ix}{5.5} ou \sga{5}{I}{3.1.2}) pour se ramener au cas
des coefficients constants (voir aussi \refext{XIII}{13-ops constant}),
on montre que le théorème est également équivalent au corollaire suivant.
(On pourrait d'ailleurs supposer l'entier $n$ premier.)

\begin{corollaire2}
\label{15-prop-devissage-classique}
Soient $X$ un schéma strictement local quasi-excellent de dimension $d$,
$U$ un ouvert affine de $X$, et $n$ un entier inversible sur $X$.
Alors,
\[
\HH^q(U,𝐙/n)=0 \text{ si } q>d.
\]
\end{corollaire2}

\subsubsection{Réduction au cas complet}
Supposons dorénavant $X$ strictement local quasi-excellent,
de complété $\chap{X}$ en le point fermé, et
fixons un ouvert affine $U$ de $X$, dont on
note $\chap{U}$ l'image inverse sur $\chap{X}$.
Le morphisme naturel de~$\widehat{X}$ dans~$X$ est régulier car~$X$ est quasi-excellent.
En appliquant le lemme de changement de base par
un morphisme régulier (\refext{XIV}{14-lemme-cb-regulier}) au diagramme cartésien
$$\xymatrix{\widehat{U} \ar[d] \ar[r] & U\ar[d] \\
\widehat{X}\ar[r] & X}$$
on obtient $\HH^q(\widehat{U},\ZZ/n\ZZ)=\HH^q({U},\ZZ/n\ZZ)$
pour tout~$q\geq 0$. Voir également \cite[7.1.1]{TUBULAR@Fujiwara}
pour une autre démonstration.

\subsubsection{}Dans les deux sections qui vont suivre, nous
allons démontrer l'énoncé \ref{15-prop-devissage-classique}
dans le cas particulier où $X$ est local nœthérien complet à
corps résiduel séparablement clos.

\section{Pureté, combinatoire des branches et descente}
\label{15-section pureté et combinatoire}

\subsection{Pureté : rappel et une application}

\subsubsection{}Nous rappelons le théorème de pureté absolue démontré par O. Gabber
(\cite{AZUMINO@Fujiwara}). Par convention, on considère le schéma vide
comme un diviseur strictement à croisements normaux dont
l'ensemble des branches est indexé par l'ensemble vide.

\begin{theoreme2}[\refext{XVI}{16-theoreme-purete-absolue}]\label{15-theoreme-purete-rappel}
Soient $X$ un schéma régulier, $Z$ un diviseur strictement à
croisements normaux de complémentaire $j:U = X\backslash Z
\hookrightarrow X$ et de branches~$\{Z_i\}_{i\in I}$, et~$n$ un entier inversible sur~$X$. Il existe des isomorphismes canoniques
\begin{eqnarray*}
\R^1j_{\star}\: (\ZZ/n\ZZ) & \isolong & \bigoplus_{i\in I} \: (\ZZ/n\ZZ)_{Z_i}(-1) \\
\R^qj_{\star}\: (\ZZ/n\ZZ) & \isolong & \bigwedge^q\: \R^1 j_{\star}\: (\ZZ/n\ZZ)
\end{eqnarray*}
\end{theoreme2}

\subsubsection{Combinatoire des branches : notations}
Soient $g:X′ → X$ un morphisme entre schémas, et $U$ un ouvert de $X$.
Notons $j:U ↪ X$ l'immersion ouverte, $j′:U′ ↪ X′$
l'immersion ouverte qui s'en déduit par changement de base
et $Z$ et $Z′$ les fermés complémentaires respectifs.
Enfin on se donne un fermé $F ⊆ Z$ dont on note $F′$
l'image inverse.

\begin{définition2}
\label{15-définition même combinatoire}
On dit que $(Z ↪ X)$ et $(Z′ ↪ X′)$ ont
\emph{même combinatoire le long de $F$}
si pour tout point géométrique $\sur{z}′$ de $F′$
d'image le point géométrique $\sur{z}$ de $F$, les propriétés suivantes sont satisfaites :
\begin{enumerate}
\item les schémas $X_{(\sur{z})}$ et $X′_{(\sur{z}′)}$ sont
réguliers ;
\item le fermé $Z_{(\sur{z})}$ est un diviseur à croisements normaux
définis par des équations $f₁,…,f_r$ ;
\item les fonctions $g^⋆f₁,…,g^⋆f_r$ forment une famille
libre de $𝔪/𝔪²$, où $𝔪$ est l'idéal maximal de $Z′_{(\sur{z}′)}$.
\end{enumerate}
\end{définition2}

\subsubsection{}Il résulte de (i) et (iii) ci-dessus
que le fermé $Z′_{(\sur{z}′)}$ est un diviseur à croisements
normaux dans $X′_{(\sur{z}′)}$, de même codimension que
$Z_{(\sur{z})}$ dans $X_{(\sur{z})}$.

\subsubsection{}
\label{15-abus terminologie}
Lorsque $X$ est un schéma sur une
base $S$, et $F$ est un fermé de ce dernier,
on s'autorise à dire « … le long de $F$ »
pour « … le long de l'image inverse $F×_S X$ ».

\begin{proposition2}
\label{15-pureté et combinatoire}
Supposons que $(Z ↪ X)$ et $(Z′ ↪ X′)$ aient
même combinatoire le long d'un fermé $F$ de $X$.
Alors, pour tout entier $n$ inversible sur $X$, le morphisme d'adjonction
\[
(\R j_⋆ 𝐙/n)_{|F′} → (\R j′_⋆ 𝐙/n)_{|F′}
\]
est un \emph{isomorphisme}.
\end{proposition2}

\begin{démo}
Quitte à localiser en des points géométriques $\sur{z}′$ et $\sur{z}$,
on peut supposer les schémas strictement locaux et le morphisme
$X′ → X$ local. Il faut alors montrer que
$\R Γ(U, 𝐙/n) → \R Γ(U′, 𝐙/n)$ est un isomorphisme. D'après
le théorème de pureté \ref{15-theoreme-purete-rappel},
il suffit de montrer que le morphisme induit sur le $\HH¹$
est un isomorphisme, ce qui résulte aussitôt de la structure de ces
groupes et de ce que la classe associée à une branche $Z_i=V(f_i)$ de $Z$
est envoyée par restriction sur la classe de la
branche $g^{-1}(Z_i)=V(g^⋆f_i)$ (cf. \refext{XVI}{16-section-gysin}).
\end{démo}

\subsection{Application du théorème de descente fléchée}
\label{15-paragraphe-descente}

\subsubsection{}
Soient $X$ un schéma strictement local nœthérien, $j:U ↪ X$
un ouvert, $i_x:x ↪ X$ l'immersion fermée, et $ε:X_∙ → X$
un hyperrecouvrement pour la topologie des altérations (\refext{II}{2-définition altL}).
La proposition suivante — où les morphismes obtenus
par changement de base sont notés de façon évidente —
est un corollaire immédiat du théorème \refext{XII}{12-cb orienté}
et du fait que la cohomologie de $U$ est la fibre en $x$ de
l'image directe par $j$.

\begin{proposition2}
\label{15-cohomologie via hyperrecouvrement}
Sous les hypothèses précédentes, le morphisme d'adjonction
\[
\R Γ(U,𝐙/n) → \R {ε_x}_⋆ \big( i_{x ∙}^⋆ \R {j_∙}_⋆ 𝐙/n \big)
\]
est un \emph{isomorphisme}.
\end{proposition2}

\subsubsection{}
Supposons maintenant donné un morphisme \emph{local} $X′ → X$ de schémas strictement
locaux nœthériens. Comme précédemment, on note $U$ un
ouvert de $X$, $Z$ son complémentaire et $x$ le point fermé de $X$.
À tout hyperrecouvrement pour la topologie des altérations
$X_∙ → X$ de $X$ est associé par changement de base un
hyperrecouvrement de $X′$.

\begin{proposition2}
\label{15-isomorphisme via HR et combinatoire}
Supposons que pour chaque entier $q ≥ 0$
les fermés $(Z_q ↪ X_q)$ et $(Z′_q ↪ X′_q)$ aient
même combinatoire le long de la fibre spéciale $(X_{q})_x$.
Alors le morphisme d'adjonction
\[
\RΓ(U,𝐙/n) → \R Γ(U′,𝐙/n)
\]
est un isomorphisme.
De plus, si l'on fait l'hypothèse précédente
pour les seuls entiers $q ≤ N+1$, où $N$ est un
entier quelconque, les morphismes
$\HH^q(U,𝐙/n) → \HH^q(U′,𝐙/n)$ sont des isomorphismes pour
chaque $q ≤ N$.
\end{proposition2}

\begin{démo}
Le premier point est conséquence immédiate
des deux propositions précédentes.
En effet, d'après \ref{15-cohomologie via hyperrecouvrement}
et l'invariance de la cohomologie par changement de corps
séparablement clos, soit
\[
\Big( \R {ε_x}_⋆ \big( i_{x ∙}^⋆ \R {j_∙}_⋆ 𝐙/n \big) \Big)_{|x′} = \R
{ε′_{x′}}_⋆ \Big( \big({i_{x∙}}^⋆ \R {j_∙}_⋆ 𝐙/n\big)_{|x′_∙}\Big),
\]
il suffit de démontrer que pour chaque $q$,
l'adjonction $(\R {j_q}_⋆ 𝐙/n)_{|{X′_q}_x} → (\R {j′_q}_⋆ 𝐙/n)_{|{X′_q}_x}$
est un isomorphisme. Cela résulte de l'hypothèse
combinatoire et de \ref{15-pureté et combinatoire}.
La variante tronquée est un corollaire de la démonstration
précédente et du fait que la cohomologie en degré $q$
ne dépend que des étages $≤q+1$.
\end{démo}

\begin{remarque2}
Dans ce critère, on ne fait d'hypothèses qu'en les points des fibres spéciales des hyper-recouvrements ; c'est ce qui en fait toute sa force.
\end{remarque2}

\section{Uniformisation et approximation des données}\label{15-section-5}

\subsection{Notations}

\subsubsection{}Soit $X,U,Z$ et $n$ comme à la fin
du ❡\ref{15-énoncés} : le schéma $X$ est strictement local
complet, $U$ en est un ouvert \emph{affine}, $Z=X-U$ (muni de la
structure réduite), et $n$ est un entier inversible sur $X$.
On veut démontrer \ref{15-prop-devissage-classique} dans ce cas.
Fixons un entier $N$.

\subsubsection{}Il résulte du théorème d'uniformisation
(\refext{VII}{7-uniformisation}), complété par
\refext{XIII}{13-cb}, qu'il existe un hyperrecouvrement
pour la topologie des altérations $ε:X_∙ → X$ tel que
chaque $X_q$ soit régulier et chaque $Z_q=Z×_X X_q$ soit
un diviseur strictement à croisements normaux. Il est
loisible de supposer les schémas $X_q$ affines.

\subsubsection{}Soient $k$ le corps résiduel de $X$,
soit $C$ un anneau de Cohen de corps résiduel $k$
(\refext{IV}{4-19.8.4}) et $S=\Spec(C)$.
Il résulte du théorème de structure des anneaux locaux
nœthériens (\egazero{iv}{19.8.8}) qu'il existe un entier $m$
et une immersion fermée de $X$ dans le complété en l'origine
(de la fibre spéciale sur $S$) de l'espace affine $𝐀^m_S$.
Notons $\chap{E}$ ce complété, $E$ son analogue hensélien
(hensélisé de l'espace affine en l'origine) et enfin $e$
le point fermé de ce dernier.

\subsubsection{}Le schéma $E$ étant quasi-excellent,
le morphisme de complétion $\chap{E} → E$ est (local) régulier
de sorte que, d'après le théorème de Popescu,
on peut écrire :
\[
\chap{E} = \lim_α E_α,
\]
où les $E_α → E$ sont des morphismes essentiellement lisses
entre schémas strictement locaux, la limite étant filtrante.
Notons que les schémas $E_α$ sont essentiellement de type
fini sur le \emph{trait} $S$.

\subsection{Passage à la limite}

\subsubsection{}
Quitte à restreindre l'ensemble d'indices, c'est-à-dire à supposer
$α ≫ 0$, les principes généraux de \ega{IV₃}{§8}, joints au
fait que les schémas $X,Z,U$ et les $X_q$ pour $q ≤ N$
sont de présentation finie sur $\chap{E}$,
entraînent l'existence de diagrammes à carrés cartésiens
\begin{eqnarray*}
\xymatrix{ U_{≤ N, α} \ar[r] \ar[d] & X_{≤ N, α}   \ar[d] & Z_{≤ N, α}\ar[l] \ar[d] \\
U_{\alpha} \ar[r] \ar[dr] & X_{\alpha} \ar[d] & Z_{\alpha} \ar[l] \ar[dl] \\
& E_α & }
\end{eqnarray*}
dont l'analogue sur $\chap{E}$ se déduit par changement de
base $\chap{E} → E_α$. De plus, on peut supposer que pour
chaque $α$, $X_α → E_α$ est une immersion fermée — de sorte que $X_α$ est strictement local —,
et $U_α → X_α$ une immersion ouverte affine de
complémentaire $Z_α$.

\begin{remarque2}
Les schémas $X_q$ et $X_{q α}$ ont même fibre spéciale pour tout~$q ≤  N$.
\end{remarque2}

\subsubsection{}Il résulte du théorème de la forme standard
(\refext{II}{2-forme standard l'-altérations})\footnote{La
démonstration dans le cas « $ℓ=1$ », qui est celui
utilisé ici, est valable sans supposer $X$ irréductible.}
et de \ega{IV₃}{8.10.5} que l'on peut supposer — quitte à
considérer de grands indices — que les $X_{≤ N, α} → X$
sont des hyperrecouvrements (tronqués) pour la topologie des
altérations.

\subsubsection{}Vérifions que l'on peut supposer
que pour chaque $α$ et chaque $q ≤ N$, le « modèle » $X_{q α}$ de $X_q$
est régulier le long de sa fibre spéciale.
Fixons $q$ puis oublions le. Le schéma $X$
est donc maintenant supposé régulier, affine, de type fini
sur $\chap{E}$. Le problème est local pour la topologie de Zariski ;
on peut donc supposer que $X$ est un fermé $V(f₁,…,f_c)$
de codimension $c$ dans un espace affine $Y=𝐀^m_{\chap{E}}$.
Quitte à supposer les $α$ suffisamment grands,
il existe des fonctions $f_{i α}$
telles que $f_i=g^⋆f_{i α}$ pour chaque $i$, où $g$
est le morphisme évident $Y → Y_α=𝐀^m_{E_α}$.
Soit $x$ un point de $Y$ appartenant à la fibre
spéciale de $X$ et soit $x_α=g(x)$ son image.
Notons $𝔪_x$ (resp. $𝔪_{x_α}$) l'idéal maximal de $𝒪_{Y,x}$
(resp. $𝒪_{Y_α,x_α}$). Par hypothèse de régularité de $X$ en $x$,
les images $\sur{f₁},…,\sur{f_c}$ des $f_i$
dans $𝔪_x/ 𝔪²_x$ sont linéairement indépendantes
sur $κ(x)=𝒪_{Y,x}/𝔪_x$. Le diagramme
\begin{center}
 \begin{tikzpicture}[auto]
  \matrix(diag)[matrix of math nodes,column sep=3em,row sep=5ex]
  {|(f)|  f_i     & |(m)|  𝔪_x     & |(t)|  𝔪_x/𝔪²_x    \\
   |(fa)| f_{i α} & |(ma)| 𝔪_{x_α} & |(ta)| 𝔪_{x_α}/𝔪²_{x_α} \\};
  \draw[->>] (m) -- (t);
  \draw[->] (ma) -- node{$g^⋆$} (m);
  \draw[->>] (ma) -- (ta);
  \draw[->] (ta) --  (t);
  \draw[|->] (fa) --  (f);
 \end{tikzpicture}
\end{center}
étant commutatif, il résulte de l'égalité $κ(x)=κ(x_α)$
que les images $\sur{f_{i α}}$ des $f_{i α}$ dans $ 𝔪_{x_α}/𝔪²_{x_α}$
sont linéairement indépendantes sur $κ(x_α)$.
Il en résulte que le sous-schéma fermé $X_α=V(f_{1 α},…,f_{c α})$
de $Y_α$ est régulier en $x_α$.


\subsubsection{}On montre de même que
l'on peut supposer que pour chaque $α$ et chaque $q ≤ N$,
les immersions $(Z_q ↪ X_q)$ et $(Z_{q α} ↪ X_{q α})$ ont même combinatoire le long
du point fermé $e ∈ E$, c'est-à-dire le long des fibres spéciales.

\subsubsection{}
Il résulte de la proposition \ref{15-isomorphisme via HR et combinatoire}
que les morphismes d'adjonction
\[
\HH^q(U_α,𝐙/n) → \HH^q(U,𝐙/n)
\]
sont des isomorphismes pour $q<N$. Nous allons montrer
que si $q>d=\dim(X)$ et $α$ est suffisamment grand,
on a $\HH^q(U_α,𝐙/n)=0$. Ceci achèvera la démonstration
du théorème de Lefschetz affine. Notons que l'ouvert
affine $U_α$ du schéma strictement local $X_α$
n'est en général pas de dimension $d$.

\subsection{Utilisation d'une section}

\subsubsection{}
Soit $σ : E → E_α$ une section de $E_α → E$
et $X^σ_α$ (resp. $U^σ_α,Z^σ_α$) le $E$-schéma déduit de $X_α$
(resp. $U_α$, $Z_α$)
par changement de base.
Notons également $X^σ_{≤ N,α}$ l'hyperrecouvrement
de $X^σ_α$ pour la topologie des altérations obtenu
à partir de $X_{≤ N,α} → X_α$ par changement de base.
On rappelle que l'on remplace les schémas obtenus
par produit fibré usuel par la réunion des adhérences
de leurs composantes irréductibles dominant une composante
irréductible de $X^σ_α$. (En d'autres termes, on applique le
foncteur $T↦ T_{\maxdom}$ de \refext{II}{2-dom}.)
Enfin $U^σ_{≤ N,α}$ (resp. $Z^σ_{≤ N,α}$) est l'ouvert (resp. le fermé)
simplicial évident.

\subsubsection{}Il résulte de \refext{III}{3-perm} et
\refext{III}{3-reduc} (démonstration)
que si $α$ est suffisamment grand et $σ$ est
suffisamment proche de l'identité,
alors, les immersions fermées
$(Z^σ_{qα} ↪ X^σ_{qα})$ et
$(Z_{qα} ↪ X_{qα})$ ont \emph{même combinatoire}
le long de la fibre spéciale au-dessus de $E$ pour
chaque $q ≤ N$.
Il en résulte comme ci-dessus que le morphisme
\[
\HH^q(U_α,𝐙/n) → \HH^q(U^σ_α,𝐙/n)
\]
est un isomorphisme. Comme l'ouvert $U^σ_α$
est affine dans $X^σ_α$ de dimension $d$
et essentiellement de type fini sur le trait $S$,
il résulte par passage à la limite
du théorème de Lefschetz (dû à O. Gabber)
\cite[2.4]{PV@Illusie}, on a
\[
\HH^q(U^σ_α,𝐙/n)=0 \text{ si } d<q<N.
\]
Finalement, $\HH^q(U,𝐙/n)=0$ si $q>d=\dim(X)$. CQFD.

\ifx\danslelivre\undefined
\bibliography{../fichiers-configuration-latex/bibliographie}
\bibliographystyle{../fichiers-configuration-latex/smfalpha}
\end{document}
\fi
